\documentclass[11pt,a4paper]{article}
\usepackage[utf8]{inputenc}
\usepackage{amsmath,amssymb,amsthm}
\usepackage{mathrsfs}
\usepackage{geometry}
\geometry{margin=1in}
\usepackage{hyperref}
\usepackage{enumitem}

\newtheorem{definition}{Definition}[section]
\newtheorem{theorem}{Theorem}[section]
\newtheorem{proposition}{Proposition}[section]
\newtheorem{remark}{Remark}[section]

\title{\textbf{Quantum Field Theory in Curved Spacetime\\and the Operator Product Expansion}\\[6pt]
\large Lecture by Robert M.\ Wald}
\author{Lecture Notes from the AQFT 50th Anniversary Conference}
\date{}

\begin{document}
\maketitle

\begin{abstract}
Robert Wald presents a comprehensive framework for quantum field theory in curved
spacetime, arguing that the algebraic approach is essential for formulating the theory
without reference to a preferred vacuum state or Hilbert space. He reviews the
difficulties of generalizing the Wightman axioms to curved spacetime, describes
how the microlocal spectrum condition and local covariance overcome these difficulties,
and proposes the operator product expansion as a fundamental replacement for the
existence of a Poincar\'e-invariant vacuum state.
\end{abstract}

\tableofcontents

%---------------------------------------------------------------------
\section{What Is QFT in Curved Spacetime?}
%---------------------------------------------------------------------

Quantum field theory in curved spacetime is a theory in which:
\begin{itemize}[nosep]
  \item All matter fields are treated fully quantum mechanically.
  \item Gravity is treated classically, as a curved spacetime metric $g_{\mu\nu}$.
\end{itemize}

This is not expected to be exact, but should be a good approximation away
from Planck scales. Remarkably, some of the deepest insights into quantum
gravity (the Hawking effect, black hole thermodynamics) emerge from this
semiclassical framework.

\subsection{The Central Difficulty}

In classical field theory, generalizing from flat to curved spacetime is
straightforward because field equations and solutions are cleanly separated:
\begin{itemize}[nosep]
  \item \textbf{Field equations} generalize locally and covariantly.
  \item \textbf{Solutions} have no natural correspondence between different spacetimes.
\end{itemize}

In quantum field theory, \emph{states play the role of solutions}, but
properties of states are deeply embedded in the formulation of the theory.
This makes generalization far more difficult.

%---------------------------------------------------------------------
\section{Problems with the Wightman Axioms in Curved Spacetime}
%---------------------------------------------------------------------

The Wightman axioms (from Streater--Wightman) are:
\begin{enumerate}
  \item States are rays in a Hilbert space carrying a representation of
        the Poincar\'e group.
  \item \textbf{Spectrum condition}: total four-momentum lies in
        $\overline{V}_+$.
  \item \textbf{Unique Poincar\'e-invariant vacuum state}.
  \item Quantum fields are operator-valued distributions on a dense,
        Poincar\'e-invariant domain.
  \item Fields transform covariantly under Poincar\'e transformations.
  \item \textbf{Local commutativity} at spacelike separations.
\end{enumerate}

\subsection{What Survives}

In a generic curved spacetime \textbf{only axiom (6)} generalizes
straightforwardly. Everything else encounters serious problems:
\begin{itemize}[nosep]
  \item \textbf{No symmetries}: generic spacetimes have no isometries,
        so Poincar\'e invariance/covariance is meaningless.
  \item \textbf{No preferred Hilbert space}: for free fields with
        non-compact Cauchy surfaces, there are unitarily inequivalent
        representations with no criterion to prefer one.
  \item \textbf{No vacuum state}: no distinguished state exists in
        the absence of symmetry.
  \item \textbf{No spectrum condition}: without time-translation symmetry,
        total energy is neither conserved nor necessarily positive.
  \item \textbf{No particle concept}: the notion of positive-frequency
        solutions depends on time-translation symmetry.
\end{itemize}

\subsection{The Spectrum Condition in Detail}

Even classically, in general relativity the conservation law
$\nabla_\mu T^{\mu\nu} = 0$ does \emph{not} yield a conserved total
energy without a Killing vector field. In quantum theory:
\begin{itemize}[nosep]
  \item The stress-energy tensor violates all pointwise energy conditions.
  \item Total energy can be negative (e.g., Casimir effect on $S^1$).
  \item No useful global positivity condition exists.
\end{itemize}

\subsection{The Vacuum State Problem}

For free fields, quasi-free Hadamard states are ``vacuum-like,'' but:
\begin{itemize}[nosep]
  \item There are many such states, none preferred.
  \item Different choices yield unitarily inequivalent representations.
  \item The situation is analogous to seeking preferred coordinates in
        general relativity---fruitless and unnecessary.
\end{itemize}

The resolution: \emph{formulate the theory algebraically}, specifying
the algebra of local field observables without choosing a state or
representation.

%---------------------------------------------------------------------
\section{Overcoming the Difficulties}
%---------------------------------------------------------------------

\subsection{The Algebraic Approach}

The Hilbert space problem is solved by simply not choosing one: formulate
the theory in terms of the algebra of observables $\mathfrak{A}$.

\subsection{The Microlocal Spectrum Condition}

The global spectrum condition is replaced by a \emph{local} condition
using microlocal analysis. The \textbf{wavefront set} $\mathrm{WF}(u)$
of a distribution $u$ captures not just where $u$ is singular but in
which \emph{cotangent directions} the singularity occurs.

For the two-point function $\omega(\Phi(x)\Phi(y))$, the microlocal
spectrum condition demands:
\[
  \mathrm{WF}\bigl(\omega(\Phi(x)\Phi(y))\bigr) \subset
  \bigl\{(x,k_x;\, y,k_y) \;\big|\; k_x \in \overline{V}_+^*\bigr\},
\]
expressing a \emph{locally positive frequency} character.

\subsection{Local Covariance}

\begin{definition}[Locally Covariant QFT]
The construction of the QFT depends on the spacetime metric only in a
\textbf{local and covariant} way: if $(M,g)$ can be isometrically and
causally embedded into $(N,g')$, then there is a natural isomorphism
of the field algebra on $M$ into that on $N$.
\end{definition}

When specialized to Minkowski spacetime with Poincar\'e transformations
as embeddings, this reproduces Poincar\'e covariance. But it is formulated
without any reference to symmetries.

%---------------------------------------------------------------------
\section{Replacing the Vacuum: The OPE Axiom}
%---------------------------------------------------------------------

\subsection{Motivation}

The existence of a Poincar\'e-invariant vacuum is used critically in:
\begin{itemize}[nosep]
  \item Spin-statistics theorem.
  \item PCT theorem.
  \item Definition of scattering theory.
\end{itemize}

An adequate replacement is needed that:
\begin{enumerate}[nosep]
  \item Makes no reference to preferred states.
  \item Captures the same physical content.
  \item Works in arbitrary globally hyperbolic spacetimes.
\end{enumerate}

\subsection{The Operator Product Expansion}

The proposal (Hollands--Wald): elevate the \textbf{operator product expansion}
(OPE) to a fundamental axiom. As points $x_1,\ldots,x_n \to y$:
\begin{equation}
  \Phi_{a_1}(x_1) \cdots \Phi_{a_n}(x_n) \sim
  \sum_c C^c_{a_1\cdots a_n}(x_1,\ldots,x_n;\, y)\, \Phi_c(y),
\end{equation}
where the $C^c_{a_1\cdots a_n}$ are \emph{state-independent} distributions
(the OPE coefficients), and $\Phi_c$ ranges over all composite fields.

\subsection{Key Properties}

The OPE coefficients must satisfy:
\begin{enumerate}[nosep]
  \item \textbf{Local covariance}: the $C$-coefficients are locally and
        covariantly constructed from the metric.
  \item \textbf{Compatibility with $*$-operation}.
  \item \textbf{Causality}: appropriate support conditions.
  \item \textbf{Scaling and dimension}: the coefficient of the identity
        in $\Phi_a(x)\Phi_a^*(y)$ defines the dimension $\dim(\Phi_a)$;
        all other coefficients are bounded in singularity by this.
  \item \textbf{Spectrum conditions} on the wavefront sets.
  \item \textbf{Associativity}: different orders of taking OPEs yield
        the same result.
\end{enumerate}

\subsection{The Role of OPE Coefficients}

The coefficients of the identity element in OPEs play the role that
\emph{vacuum expectation values} play in Minkowski-space QFT. The algebra
itself is quite ``trivial''---it is the free $*$-algebra generated by the
fields, factored only by:
\begin{itemize}[nosep]
  \item Relations derived from the OPE (e.g., field equations).
  \item Commutativity/anti-commutativity at spacelike separations.
\end{itemize}

All nontrivial information about the theory is contained in the OPE
coefficients.

%---------------------------------------------------------------------
\section{Existence and Uniqueness}
%---------------------------------------------------------------------

\subsection{Existence}

The OPE is known to exist:
\begin{itemize}[nosep]
  \item For free fields in curved spacetime (exact).
  \item Order by order in perturbation theory for renormalizable interacting
        fields (Hollands).
\end{itemize}

It is \emph{not} expected to exist for non-renormalizable theories.

\subsection{Uniqueness}

\begin{theorem}[Hollands--Wald]
For a given field content and set of OPE coefficients satisfying all
axioms, the state space (positive linear functionals satisfying the
microlocal spectrum condition and OPE relations) is nonempty. If two
sets of OPE coefficients agree in one spacetime, they agree in all
spacetimes (by local covariance and the time-slice axiom).
\end{theorem}

%---------------------------------------------------------------------
\section{Spin-Statistics and PCT on Curved Spacetimes}
%---------------------------------------------------------------------

Using the OPE framework:

\begin{theorem}[Spin-Statistics, Verch]
If a locally covariant QFT satisfies the Wightman axioms in Minkowski
spacetime, it has the correct spin-statistics connection in every
globally hyperbolic spacetime.
\end{theorem}

\begin{theorem}[PCT, Hollands]
Given an OPE around each spacetime point, there exists an anti-linear
map relating the OPE of the field to that of the charge-conjugate field
on the same spacetime with reversed spacetime orientation.
\end{theorem}

%---------------------------------------------------------------------
\section{Summary: The Modern Viewpoint}
%---------------------------------------------------------------------

\begin{center}
\begin{tabular}{l|l}
\textbf{Wightman (Minkowski)} & \textbf{Hollands--Wald (Curved)} \\
\hline
Hilbert space + Poincar\'e group & Algebraic approach + local covariance \\
Spectrum condition & Microlocal spectrum condition \\
Unique vacuum state & OPE axiom \\
Poincar\'e covariance of fields & Locally covariant fields \\
Dense Poincar\'e-invariant domain & GNS construction from any state \\
Local commutativity & Local commutativity (unchanged) \\
\end{tabular}
\end{center}

The operator product expansion, elevated to a fundamental axiom, provides
a complete replacement for the vacuum state, capturing all the physical
content of quantum field theory in a framework that works in arbitrary
globally hyperbolic spacetimes.

\end{document}
